\documentclass[12pt,reqno]{amsart}

\usepackage{amsmath,amssymb,amsthm}
\usepackage{booktabs}
\usepackage[colorlinks=true,linkcolor=blue,citecolor=blue,urlcolor=blue]{hyperref}
\usepackage[margin=1.15in]{geometry}
\usepackage{enumitem}
\usepackage{microtype}

\newtheorem{theorem}{Theorem}[section]
\newtheorem{proposition}[theorem]{Proposition}
\newtheorem{corollary}[theorem]{Corollary}
\theoremstyle{definition}
\newtheorem{definition}[theorem]{Definition}
\theoremstyle{remark}
\newtheorem{remark}[theorem]{Remark}
\newtheorem{question}[theorem]{Question}

\newcommand{\ZZ}{\mathbb{Z}}
\newcommand{\QQ}{\mathbb{Q}}
\newcommand{\Gstar}{G^{\!*}}
\newcommand{\varpiup}{\varpi}
\newcommand{\agm}{\operatorname{agm}}
\newcommand{\abs}[1]{\lvert #1 \rvert}

\title[The Missing Ratio]{The Missing Ratio:\\[2pt]
  $\Gamma(\tfrac14)/\Gamma(\tfrac34)$ as a Fundamental Constant}

\author{}
\address{}

\subjclass[2020]{33B15, 11M06, 01A55, 33E05}
\keywords{Gamma function, reflection formula, lemniscate constant,
  Euler, fundamental constants, Wallis product}

\date{\today}

\begin{document}

\begin{abstract}
The Euler reflection formula at $z=\tfrac14$ gives
$\Gamma(\tfrac14)\,\Gamma(\tfrac34) = \pi\sqrt{2}$.
The associated classical \emph{ratio}
$\Gamma(\tfrac14)/\Gamma(\tfrac34) \approx 2.9587$---which
we denote $\Gstar$ and call the \emph{lemniscatic bridge constant}---admits
a Wallis-type infinite product, appears as a scaled period of the CM elliptic curve
$y^{2}=x^{3}-x$, and factors the lemniscate constant as
$\varpiup = {\Gstar}\sqrt{\pi}/2$.

We introduce this normalisation alongside $\pi$ and $\varpiup$ through
its classical identities. The dedicated notation is a pedagogical
proposal; no novelty, priority, or historical-neglect claim is made.
\end{abstract}

\maketitle
\tableofcontents


% ================================================================
\section{The reflection formula produces two objects}
% ================================================================

The Euler reflection formula states that for all $z\notin\ZZ$,
\begin{equation}\label{eq:reflection}
  \Gamma(z)\,\Gamma(1-z) \;=\; \frac{\pi}{\sin(\pi z)}\,.
\end{equation}
At $z = \tfrac14$:
\begin{equation}\label{eq:product}
  \Gamma(\tfrac14)\,\Gamma(\tfrac34)
  \;=\; \frac{\pi}{\sin(\pi/4)}
  \;=\; \pi\sqrt{2}\,.
\end{equation}
This is the \emph{product}. It yields a circular quantity ($\pi$,
modified by $\sqrt{2}$).

The same two Gamma values also determine a \emph{ratio}:
\begin{equation}\label{eq:ratio}
  \frac{\Gamma(\tfrac14)}{\Gamma(\tfrac34)}
  \;=\; \frac{\Gamma(\tfrac14)^{2}}{\pi\sqrt{2}}
  \;=\; 2.9586751191886388\ldots
\end{equation}
We use $\Gstar$ as a dedicated symbol for this ratio, keeping it
distinct from both the circular constant and the lemniscatic period.

\begin{remark}
Every evaluation of the reflection formula at a rational argument
$z = p/q$ produces a product and a ratio.  At $z = \tfrac12$, the
product $\Gamma(\tfrac12)^{2} = \pi$ yields the circle constant itself,
and the ratio $\Gamma(\tfrac12)/\Gamma(\tfrac12) = 1$ is trivial.
At $z = \tfrac14$, the product yields $\pi\sqrt{2}$ and the ratio
yields $\Gstar$.  This is the \emph{first} non-trivial ratio produced
by the reflection formula at a point of the form $z = 1/(2n)$.
\end{remark}


% ================================================================
\section{Periods and normalisations}
% ================================================================

Three complementary descriptions make the relationship between the
ratio, the period, and the circular constant explicit.

\subsection{The geometric period}

The lemniscate constant is a geometric period:
\begin{equation}\label{eq:varpi}
  \varpiup \;=\; \frac{\Gamma(\tfrac14)^{2}}{2\sqrt{2}\,\Gamma(\tfrac12)}
  \;=\; 2.6220575542921198\ldots
\end{equation}
is the half-perimeter of the lemniscate $r^{2}=\cos 2\theta$ and a
period of the elliptic curve $y^{2}=x^{3}-x$.  As a period, it has
geometric meaning: it is a length one can draw. The constant $\Gstar$
is the period rescaled by $2/\sqrt\pi$, as well as a ratio of two
Gamma values.

\subsection{The reflection product}

The product $\Gamma(\tfrac14)\Gamma(\tfrac34)=\pi\sqrt2$ fixes the
product of the two Gamma values. Their ratio records their relative
size. Both descriptions are useful, and neither alone determines
the two values without further information.

\subsection{Making the normalisation visible}

The equivalent forms
$\varpiup=\Gamma(\tfrac14)^2/(2\sqrt{2\pi})$ and
$\varpiup=\Gstar\sqrt\pi/2$ emphasise different relations.
The second makes the chosen Gamma-ratio normalisation explicit.


% ================================================================
\section{The constant and its properties}
% ================================================================

\begin{definition}\label{def:Gstar}
The \emph{lemniscatic bridge constant} is
\begin{equation}\label{eq:Gstar}
  {\Gstar} \;=\; \frac{\Gamma(\tfrac14)}{\Gamma(\tfrac34)}
  \;=\; 2.9586751191886388\ldots
\end{equation}
\end{definition}

We collect its principal representations, all classical:
\begin{align}
  {\Gstar} &= \frac{\Gamma(\tfrac14)^{2}}{\sqrt{2}\,\Gamma(\tfrac12)^{2}}
    &&\text{($\pi$-free form)}\label{eq:pifree}\\[4pt]
  {\Gstar} &= \frac{2\varpiup}{\sqrt{\pi}}
    &&\text{(lemniscate form)}\label{eq:lem}\\[4pt]
  {\Gstar} &= \frac{2\sqrt{2}}{\sqrt{\pi}}\,
    K\!\bigl(\tfrac{1}{\sqrt{2}}\bigr)
    &&\text{(elliptic integral)}\label{eq:elliptic}\\[4pt]
  {\Gstar} &= \frac{2\sqrt{\pi}}{\agm(1,\sqrt{2})}
    &&\text{(AGM)}\label{eq:agm}\\[4pt]
  {\Gstar} &= \sqrt{2\pi}\;\theta_{3}(e^{-\pi})^{2}
    &&\text{(theta function at self-dual nome)}\label{eq:theta}\\[4pt]
  {\Gstar} &= \frac{8\,L(E,1)}{\sqrt{\pi}}
    &&\text{($L$-function of $y^{2}=x^{3}-x$)}\label{eq:Lfunction}
\end{align}

These are classical representations of the normalisation used here.
For $E\colon y^2=x^3-x$ over $\QQ$, the exact evaluation
$L(E,1)=B(\frac12,\frac14)/8=\varpiup/4$
follows from the Mellin transform of $\eta(4\tau)^2\eta(8\tau)^2$
\cite[Lemma~3.1]{LiLongTu}; no BSD assumption is needed for this identity.


% ================================================================
\section{What the ratio knows that the product does not}
% ================================================================

The product $\Gamma(\tfrac14)\Gamma(\tfrac34)$ is \emph{symmetric}: it
treats $z=\tfrac14$ and $z=\tfrac34$ on equal footing.  Swapping them
changes nothing. The product yields $\pi\sqrt2$ at these two arguments.
For general $z$, it still depends on $z$ through $\pi/\sin(\pi z)$;
the fixed sum $z+(1-z)=1$ does not determine the product.

The ratio $\Gamma(\tfrac14)/\Gamma(\tfrac34)$ is \emph{asymmetric}: it
distinguishes $z=\tfrac14$ from $z=\tfrac34$.  It measures \emph{how much
larger} $\Gamma(\tfrac14)$ is than $\Gamma(\tfrac34)$.  This asymmetry
carries information:

\begin{theorem}[Wallis product for the ratio]\label{thm:wallis}
\begin{equation}\label{eq:wallis-ratio}
  {\Gstar} \;=\; \lim_{N\to\infty}\;(N+1)^{-1/2}
  \prod_{k=0}^{N}\frac{4k+3}{4k+1}\,.
\end{equation}
\end{theorem}

\begin{proof}
Write $Q_{N} = \prod_{k=0}^{N}(4k+3)/(4k+1)
= (\tfrac34)_{N+1}/(\tfrac14)_{N+1}
= \Gamma(\tfrac34+N+1)\Gamma(\tfrac14)/
[\Gamma(\tfrac14+N+1)\Gamma(\tfrac34)]$.
By the asymptotic ratio $\Gamma(x+a)/\Gamma(x+b)\sim x^{a-b}$,
$Q_{N}\sim [\Gamma(\tfrac14)/\Gamma(\tfrac34)]\,(N+1)^{1/2}$.
\end{proof}

The numerators $3, 7, 11, 15, 19, 23, \ldots$ are integers
$\equiv 3\pmod{4}$.  The denominators $1, 5, 9, 13, 17, 21, \ldots$
are integers $\equiv 1\pmod{4}$. For an odd \emph{prime} $p$, these
residue classes distinguish inert primes ($p\equiv3\pmod4$) from
split primes ($p\equiv1\pmod4$) in $\ZZ[i]$. The displayed product,
however, runs over all integers in these progressions, including
composites. It is a residue-class product, not an Euler product over
primes or a measure of an advantage of inertness over splitting.


% ================================================================
\section{The triad and the factorisation of \texorpdfstring{$\varpiup$}{varpi}}
% ================================================================

\begin{theorem}[Factorisation of $\varpiup$]\label{thm:factor}
\begin{equation}\label{eq:factor}
  \varpiup \;=\; \frac{{\Gstar}\cdot\sqrt{\pi}}{2}\,.
\end{equation}
\end{theorem}

\begin{proof}
$\Gstar\sqrt{\pi}/2
= [\Gamma(\tfrac14)/\Gamma(\tfrac34)]\cdot\Gamma(\tfrac12)/2$.
By reflection, $\Gamma(\tfrac34)=\pi\sqrt{2}/\Gamma(\tfrac14)$, so
$\Gstar\sqrt{\pi}/2
= \Gamma(\tfrac14)^{2}\Gamma(\tfrac12)/(2\pi\sqrt{2})
= \Gamma(\tfrac14)^{2}/(2\sqrt{2}\Gamma(\tfrac12)) = \varpiup$.
\end{proof}

This algebraic factorisation uses two residue-class product limits:
\begin{itemize}
  \item $\sqrt{\pi}$, which arises from the Wallis product over the
    even/odd partition (the mod-$2$ ``race'');
  \item $\Gstar$, which arises from the Wallis product over the
    $1\bmod 4$ / $3\bmod 4$ partition (the mod-$4$ ``race'').
\end{itemize}
Multiplying the two limits and dividing by two gives $\varpiup$.
The factorisation alone establishes no independence statement.

\begin{proposition}[A finite arithmetic-progression product]
\label{prop:no-third}
For $a,b,c>0$ and an integer $N\geq1$, let
$P_N=\prod_{n=0}^{N-1}(an+b)/(an+c)$. Then
\[
 P_N=\frac{\Gamma(N+b/a)\Gamma(c/a)}
 {\Gamma(b/a)\Gamma(N+c/a)},\qquad
 \lim_{N\to\infty}N^{(c-b)/a}P_N
 =\frac{\Gamma(c/a)}{\Gamma(b/a)}.
\]
\end{proposition}

\begin{proof}
Cancel $a$ in each factor and use
$\prod_{n=0}^{N-1}(n+u)=\Gamma(N+u)/\Gamma(u)$.
The limit follows from
$\Gamma(N+u)/\Gamma(N+v)\sim N^{u-v}$.
\end{proof}
This reduction does not classify all special values of Gamma ratios
and gives no nonexistence theorem for representations of $\varpiup$.


% ================================================================
\section{The triad identity}
% ================================================================

\begin{theorem}\label{thm:triad}
$\pi = 4\varpiup^{2}/{\Gstar}^{2}$.
\end{theorem}

\begin{proof}
$4\varpiup^{2}/{\Gstar}^{2}
= 4\cdot({\Gstar}\sqrt{\pi}/2)^{2}/{\Gstar}^{2}
= 4\cdot{\Gstar}^{2}\pi/(4{\Gstar}^{2}) = \pi$.
\end{proof}

This expresses the circle constant in terms of two lemniscatic
quantities. It rearranges the same period normalisation; it does not
establish an independent order in which the constants are derived.


% ================================================================
\section{Why this matters}
% ================================================================

The constant $\Gstar = \Gamma(\tfrac14)/\Gamma(\tfrac34)$ appears
identically in at least seven branches of mathematics (see~\cite{Paper2}
for a catalogue of fifty-two representations):
\begin{center}
\begin{tabular}{ll}
\toprule
\textbf{Branch} & \textbf{Form} \\
\midrule
Complex analysis & Gamma ratio $\Gamma(\tfrac14)/\Gamma(\tfrac34)$ \\
Classical geometry & Scaled lemniscate arc length $2\varpiup/\sqrt{\pi}$ \\
Elliptic integrals & $2\sqrt{2}\,K(1/\!\sqrt{2})/\sqrt{\pi}$ \\
Computational number theory & $2\sqrt{\pi}/\agm(1,\sqrt{2})$ \\
Modular forms & $\sqrt{2\pi}\,\theta_{3}(e^{-\pi})^{2}$ \\
Lattice combinatorics & $\sqrt{2\pi\,W_{3}}$ (Watson integral) \\
Arithmetic geometry & $8\,L(E,1)/\sqrt{\pi}$ (CM $L$-value) \\
\bottomrule
\end{tabular}
\end{center}

The product and ratio of the same two Gamma values organise these
connections in different ways. The period normalisation, the
residue-class product, and the specified CM-curve $L$-value each
offer a useful route to the same constant.

Calling $\Gstar$ fundamental is a judgement about mathematical
usefulness, not a further theorem. The classical identities collected
here and in the companion papers~\cite{Paper2,Paper3} support a
pedagogical case for giving this normalisation a dedicated symbol.


% ================================================================
\section{A note on naming}
% ================================================================

We have used the notation $\Gstar$ throughout, following its role as the
``bridge constant'' in Foundational Ternary Dynamics.  We note that the
constant admits a simpler description than any notation suggests:

\medskip
\begin{center}
\emph{$\Gstar$ is the ratio of the Gamma function at the
quarter-point to its reflection.}
\end{center}
\medskip

\noindent
A dedicated symbol can make the relationships among these classical
representations easier to read. We propose it for that purpose,
without claiming a new constant or established naming convention.


% ================================================================
% REFERENCES
% ================================================================

\begin{thebibliography}{10}

\bibitem{BorweinBorwein}
J.~M.~Borwein and P.~B.~Borwein,
\emph{Pi and the AGM}, Wiley, 1987.

\bibitem{Finch}
S.~R.~Finch,
\emph{Mathematical Constants}, Cambridge Univ.\ Press, 2003.

\bibitem{LiLongTu}
W.-C.~W.~Li, L.~Long, and F.-T.~Tu,
\emph{Computing Special $L$-Values of Certain Modular Forms with Complex
Multiplication}, SIGMA \textbf{14} (2018), 090, Lemma~3.1.
\url{https://doi.org/10.3842/SIGMA.2018.090}.

\bibitem{Paper2}
\emph{Fifty-two faces of the lemniscatic bridge constant},
companion paper.

\bibitem{Paper3}
\emph{The two races: Wallis products for $\sqrt{\pi}$, $\Gstar$, and
the composite nature of $\varpiup$},
companion paper.

\bibitem{Wallis}
J.~Wallis,
\emph{Arithmetica Infinitorum}, Oxford, 1656.

\end{thebibliography}

\end{document}
